\scanpage{313}
\mspage{12}
Alors la connaissance de $C_{\cdot\deleted{\unsure{*}}}=\mathbb C_{h\cdot}(X_*)
=\mathbb C_\cdot\otimes_{\mathbb Z}h$ est encore suffisante pour
la construction de $F(X_*)$, car on aura
évidemment
\[
\begin{aligned}
F(X_*)&\simeq\operatorname{Hom}_{h\cdot}(C_{\cdot h},K_\cdot)\\
&\simeq Z^0\bigl(\underbrace{\mathcal Hom^{\cdot}_{h\cdot}(C_{\cdot h},K_\cdot)}_{
\substack{\mathcal Hom^i_{hj}(C_{\cdot h}^{\deleted{\unsure{i}}},K_\cdot)\\
=\operatorname{Hom}_h(C_{ih},K_j)}}\bigr)
\end{aligned}
\]
($\operatorname{Hom}_h(A,M)$, pour $A\in\mathrm{Ab}_h$, et $M\in B$, étant
défini par
\[
\operatorname{Hom}_B(N,\operatorname{Hom}_h(A,M))\simeq
\operatorname{Hom}_h(A,\operatorname{Hom}_B(N,M)))\,.
\]

$2^\circ)$ Si on regarde toutes les façons possibles
\deleted{\unsure{d'envisager}}\inserted{d'envisager}
\[
(\mathrm{Ss})^o\xrightarrow{F}B
\]
dans $B$ stable par $\varprojlim$, avec $F$ y
commutant, il y en a \deleted{\unsure{bien sûr}} une
``universelle'', dont toutes les autres sont
``spécialisées'', obtenue en prenant
\[
B_o=(\mathrm{Ab}_\cdot)^o,\quad\text{et}\quad
F_o:(\mathrm{Ss})^o\longrightarrow B_o=(\mathrm{Ab}_\cdot)^o
\]
défini par
\[
F_o(X_*)=\mathbb C_\cdot(X_*)
\]
(qui dépend de $X_*$ de façon contravariante
quand $\mathbb C_\cdot(X_*)$ s'interprète comme
objet de $(\mathrm{Ab}_\cdot)^o$). De façon précise :

\uline{Prop. 3} Considérons
\[
\mathrm{Ss}^o\xrightarrow{F_o}B_o=(\mathrm{Ab}_\cdot)^o\simeq(\mathrm{Ab}^o)^\cdot
\]
défini par $F_o(X_*)=\mathbb C_\cdot(X_*)$. Alors
$B_o$ est une catégorie \inserted{add} avec $\varprojlim$, et
pour tte autre telle catégorie $B$,
le foncteur
\[
\underline{\operatorname{Hom}}_{\substack{\mathrm{comm.}\\\text{aux }\varprojlim}}
(\deleted{\unsure{\mathrm{Ss}}}B_o,B)\longrightarrow
\underline{\operatorname{Hom}}_{\mathrm{comm}\,\varprojlim}(\mathrm{Ss}^o,B)
\]
\margin{$\unsure{\simeq|21}$ $(B_\cdot\text{ par prop. 2})$}
donné par $\varphi\longmapsto\varphi\circ F_o$
est une équivalence de catégories.
\margin{\unsure{pourrais passer par} l'analogue avant la théorie de \uline{Dold-Puppe}}

\deleted{\unsure{Compte tenu de ceci}}
\deleted{Compte tenu de prop. 2, cela signifie
\unsure{le fait déjà que le}
(Cela signifie que $F_o$ établit l'\unsure{is}
Grâce \unsure{donc} à l'équivalence entre les deux
\unsure{cat. considérées} (\unsure{bien l'indice ! indique la commut.}
des fonct. aux $\varprojlim$)
\unsure{Hom}}
\[
\unsure{\boxed{\underline{\operatorname{Hom}}_!(B_o,B)}}
\]
Grâce au th. de Dold-Puppe, on peut
remplacer dans la déf. de
$F_o,B_o$, $(\mathrm{Ab}_\cdot)^o$ par $(\mathrm{Ab}_*)^o=\widehat{\boldsymbol\Delta}_{ab}^{\,o}$,
$\mathbb C_\cdot(X_*)$ par $\mathbb C_*(X_*)\inserted{=\mathbb Z(X_*)}$, et l'équivalence à établir
devient
\[
\underline{\operatorname{Hom}}^{!}_{\deleted{\unsure{!|\mathbb Z}}}
(\widehat{\boldsymbol\Delta}_{ab}^{\,o},B)
\xrightarrow{\overset{?}{\simeq}}
\underline{\operatorname{Hom}}^{!}_{\deleted{\unsure{!|\mathbb Z}}}(\widehat{\boldsymbol\Delta}^{\,o},B)
\]
où l'exposant ! signifie qu'on se borne
aux foncteurs qui commutent aux $\varprojlim$.
Or sous cette forme, l'assertion est
valable quelle que soit la cat. (\unsure{ess. petite})
$\boldsymbol\Delta$, \deleted{\unsure{la démonstration se faisant comme}}
\inserted{car le foncteur ``\uline{abélianisation}''}
\[
X\longrightarrow\mathbb Z(X_\cdot)\qquad
\widehat{\boldsymbol\Delta}\longrightarrow\widehat{\boldsymbol\Delta}_{ab}
\]
\deleted{\unsure{comme d'ailleurs}}
\inserted{commutant aux $\varinjlim$}
\deleted{\unsure{de} $\widehat{\boldsymbol\Delta}^{\,o}\longrightarrow B=(\widehat{\boldsymbol\Delta}_{ab})^o$}
\margin{314}
